\section{Conclusions}
\label{sec:conclusions}
Let us  discuss some perspectives for future work. The most natural extension is to consider regular functions. From a transducer perspective, this is a very important class (see the lecture notes~\cite[Part C]{Bojanczyk2026} for an overview of its numerous characterisations). The corresponding compositional syntax of \emph{regular terms}, which was introduced in~\cite{BojanczykDaviaudKrishna2018}, can be obtained from the syntax in \cref{sec:syntax} by adding two new \kl{atomic terms}: list reversal of type $A^* \to A^*$, and a diagonal map $A \to A \times A$.  We would like to extend our equational axiomatisation to cover the regular terms. Since the equivalence problem for regular functions is decidable~\cite[Theorem 1]{Gurari1982}, we see no fundamental reason why such an axiomatisation should not exist, although we expect that designing it and proving its completeness will require substantial new ideas.

A more ambitious goal would be to consider the polyregular functions~\cite{Bojanczyk2022}, which are a transducer class with polynomial output size, as opposed to the linear output size for the rational and regular functions. The corresponding compositional syntax is obtained by adding, on top of the regular terms, one more atomic term of type 
$
A^* \to (A^* \times A^*)^*,
$
which produces all possible partitions of the input list into two parts: a prefix and a suffix. This goal is much more ambitious, since the equivalence problem for polyregular functions is not known to be decidable, and thus finding a (well-behaved) equational axiomatisation would need to solve this difficult problem. However, we believe that the syntactic methods and ideas about canonical models that we developed in this paper could  be useful for this goal, and might be a new line of attack on problems where the automata method for transducer equivalence has stalled. A similar situation arises for data structures beyond lists, such as trees or graphs, where the automata method has been applied only with limited success~\cite{BojanczykSchmude2020}. 

The completeness of the theory proved in this paper is about ground terms, i.e.~terms without free variables. One could also ask about equalities that involve free variables, which range over \kl{rational functions}. We do not know about completeness for such a theory, or even decidability. For more expressive variants, in which variables also range over functions but more operations are allowed, the theory can become undecidable~\cite[Theorem 3.1]{BojanczykFabinaskiStefanski2026}.

Another interesting problem is to describe the models of our axiomatisation. Kleene algebra admits many computationally relevant models besides the regular languages, such as the algebra of binary relations on a set and the algebra of square matrices over a Kleene algebra. We do not know whether our axiomatisation has similarly interesting models beyond the rational functions, and we leave this question as a direction for future investigation.

Finally, we state an open problem that is not about the equational axiomatisation of transducers, but about the simpler setting of regular expressions and Kleene algebra. There is a star-free variant of regular expressions, in which the Kleene star is not allowed, but complementation is allowed. This variant has the same expressive power as first-order logic~\cite[Theorem 10.5]{McNaughtonPapert71}. Is there a finite equational axiomatisation of star-free regular expressions? It is unclear to us whether the methods of this paper could be of any use in resolving this problem.